%%
%% fall-lecture16.tex
%% 
%% Made by Alex Nelson <pqnelson@gmail.com>
%% Login   <alex@lisp>
%% 
%% Started on  2026-11-03T10:34:10-0800
%% Last update 2026-11-03T10:34:10-0800
%% 

\lecture[Topological Vector Spaces]{}

\begin{node}
We have been looking at normed vector spaces. We begin looking at
their generalizations. The most generalized thing is a topological
vector space. A natural generalization of normed vector spaces are
``multi-normed (or poly-normed) vector spaces''. (I only managed to
find Szilasi and Lovas~\cite[p.1077]{szilasi2007some} discuss ``multinormed
vector spaces'' in this sense.)
\end{node}

\begin{definition}
A \define{Multi-Normal Vector Space} $(X,\{p_{\lambda}\}_{\lambda\in\Lambda})$
where $X$ is a vector space and $p_{\lambda}$ is an
infinite\footnote{Finite families are not interesting because we can
combine them together into a single seminorm.} family of
seminorms. We equip $X$ with the weakest topology making every
$p_{\lambda}$ continuous.

We can construct a neighborhood base of zero,
\begin{equation}
N_{\varepsilon}=\{x\in X\mid\forall\lambda\in\Lambda\ldotp p_{\lambda}(x)<\varepsilon\},
\end{equation}
or for $\lambda\in\Lambda$,
\begin{equation}
N_{\varepsilon,\lambda}=\{x\in X\mid p_{\lambda}(x)<\varepsilon\}
\end{equation}
We can describe this topology in terms of nets: $x_{\alpha}\to x$
iff for every seminorm $\forall\lambda\in\Lambda\ldotp p_{\lambda}(x_{\alpha}-x)\to0$.
We see $N_{\varepsilon,\lambda}$ are just generalized [open] balls which are
convex. Convexity for one seminorm $N_{\varepsilon,\lambda}$ (follows by triangle inequality),
and for several seminorms it amounts to the intersection of several
convex sets (which is again convex).
\end{definition}

\begin{fact}
If $X$ is a topological vector space such that every point $x\in X$
has a convex neighborhood base [i.e., a neighborhood base of convex
  neighborhoods], then there exists a family of seminorms $\{p_{\lambda}\}_{\lambda\in\Lambda}$
that generate the topology.
\end{fact}

\begin{remark}
For this reason, \define{Locally Convex [topological] Vector Space} (LCVS) is used
instead of ``multinormed vector spaces''.

[Alternatively, one can define: a topological vector space $X$
is \define{Locally Convex} if there is a base for the topology
consisting of convex sets, that is, sets $A$ such that if $x,y\in A$
then $tx+(1-t)y\in A$ for all $0<t<1$. Most topological vector spaces
that arise in practice are locally convex and Hausdorff.]
\end{remark}

\begin{framed}
Some more notes to myself, which might be useful.

\begin{definition}
Let $X$ be a vector space over $\FF$ (where $\FF$ is usually either
$\RR$ or $\CC$---we only require the existence of an absolute value
$|-|\colon\FF\to\RR$ satisfying the obvious properties, but it is also
convenient to insist $\FF$ is a topological field).
A \define{Seminorm} on $X$ is a function $\|-\|\colon X\to[0,\infty)$
such that
\begin{enumerate}
\item for all $x,y\in X$ we have $\|x+y\|\leq\|x\|+\|y\|$
\item for all $x\in X$ and $\lambda\in\FF$ we have $\|\lambda x\|=|\lambda|\cdot\|x\|$
[which implies $\|0\|=0$]
\end{enumerate}
A seminorm with the additional property $\forall x\in X\ldotp \|x\|=0\implies x=0$
is a norm.
\end{definition}

\begin{node}
Also note: we could have defined $\|-\|\colon X\to\RR$ and deduced
  $\forall x\in X\ldotp\|x\|\geq0$. This is because
  $\|-x\|=|-1|\|x\|=\|x\|$ and
  $0=\|0\|=\|x-x\|\leq\|x\|+\|-x\|=2\|x\|$ so provided $2\neq0$ in
  $\FF$ we can conclude $0\leq\|x\|$ by dividing through by $2$.
\end{node}

\begin{theorem}[{Folland~\cite[Prop~5.14]{folland1999real}}]
Let $X$ be a vector space. Let $\{p_{\alpha}\}_{\alpha\in A}$ be a
family of seminorms on $X$. Let us write the family of sets
\begin{equation}
U_{x,\alpha,\varepsilon}:=\{y\in X\mid p_{\alpha}(y-x)<\varepsilon\}
\end{equation}
where $x\in X$, $\alpha\in A$, and $\varepsilon>0$. Let $\topology$ be
the topology generated by the sets $U_{x,\alpha,\varepsilon}$. Then
the following all hold:
\begin{enumerate}
\item For each $x\in X$, the finite intersections of the sets
  $U_{x,\alpha,\varepsilon}$ (where $\alpha\in A$ and $\varepsilon>0$)
  form a neighborhood base at $x$.
\item For any net $x_{\lambda}$ in $X$, we have $x_{\lambda}\to x$ iff
  for all $\alpha\in A$ we have $p_{\alpha}(x_{\lambda}-x)\to0$.
\item We have $(X,\topology)$ forms a locally convex topological
  vector space.
\end{enumerate}
\end{theorem}

\begin{proof}
\begin{enumerate}
\item\label{step:pf:fall:lec16:math110a:seminorms-give-convex-space:1}
  If $x\in\bigcap^{k}_{j=1}U_{x_{j},\alpha_{j},\varepsilon_{j}}$, let $\delta_{j}=\varepsilon_{j}-p_{\alpha_{j}}(x-x_{j})$.
  By triangle inequality, we have $x\in\bigcap^{k}_{j=1}U_{x,\alpha_{j},\delta_{j}}\subset\bigcap^{k}_{j=1}U_{x_{j},\alpha_{j},\varepsilon_{j}}$.
  (Recall, for any family $\mathcal{E}\subset\powerset{X}$, the
  topology $\topology(\mathcal{E})$ generated by $\mathcal{E}$
  consists of $\emptyset$, $X$, and all unions of finite intersections
  of elements of $\mathcal{E}$. From this fact, we get the result.)
\item\label{step:pf:fall:lec16:math110a:seminorms-give-convex-space:2} From the previous step~\ref{step:pf:fall:lec16:math110a:seminorms-give-convex-space:1},
it suffices to observe $p_{\alpha}(x_{\lambda}-x)\to0$ if and only if
$x_{\lambda}$ is eventually in $U_{x,\alpha,\varepsilon}$ for every $\varepsilon>0$.
\item \textsc{Claim 1: it's a topological vector space}. We need to prove
  vector addition and scalar multiplication are continuous.
  Continuity of the vector operations follows easily from
  part~\ref{step:pf:fall:lec16:math110a:seminorms-give-convex-space:2}.
  Just look at nets $x_{\lambda}\to x$ and $y_{\lambda}\to y$, and
  then look at $p_{\alpha}((x_{\lambda}+y_{\lambda})-(x+y))\leq p_{\alpha}(x_{\lambda}-)+p_{\alpha}(y-y_{\lambda})\to0$
  which proves $x_{\lambda}+y_{\lambda}\to x+y$. Similarly, if
  $c_{\lambda}\to c$ is a net of scalars which converges, then
  eventually $|c_{\lambda}|\leq|c|+1$, so
  $p_{\alpha}(c_{\lambda}x_{\lambda}-cx)\leq(|c|+1)p_{\alpha}(x_{\lambda}-x)+|c_{\lambda}-c|p_{\alpha}(x)$
  which implies $c_{\lambda}x_{\lambda}\to cx$.
  
  \textsc{Claim 2: local convexity}.
  The sets $U_{x,\alpha,\varepsilon}$ are convex, so local convexity of
the topology follows from part~\ref{step:pf:fall:lec16:math110a:seminorms-give-convex-space:1}.\qedhere
\end{enumerate}
\end{proof}

Remember how linear operators $T\colon X\to Y$ between topological
vector spaces $X$ and $Y$ are continuous if and only if $T$ is bounded
(if and only if $T$ is continuous at $0$)? Well, we have a fun
analogous proposition.

\begin{proposition}[{Folland~\cite[Prop~5.15]{folland1999real}}]
Let $X$ be a topological vector space whose topology is defined by
seminorms $\{p_{\alpha}\}_{\alpha\in A}$.
Let $Y$ be a topological vector space whose topology is defined by
seminorms $\{q_{\beta}\}_{\beta\in B}$.
Let $T\colon X\to Y$ be a linear map.
Then $T$ is continuous if and only if for each $\beta\in B$ there
exists $\alpha_{1},\dots,\alpha_{k}\in A$ and $C>0$ such that [for
each $x\in X$ we have]
$q_{\beta}(Tx)\leq C\sum^{k}_{j=1}p_{\alpha_{j}}(x)$.
\end{proposition}
\end{framed}

\begin{fact}[{Folland~\cite[Prop~5.16(a)]{folland1999real}}]
  Let $X$ be a locally convex vector space.
  Then $X$ is Hausdorff if and only if for every nonzero vector $x\in X$
  there exists  a seminorm $p_{\lambda}$ such that $p_{\lambda}(x)\neq0$.
\end{fact}

\begin{fact}[{Folland~\cite[Prop~5.16(b)]{folland1999real}}]
If $X$ is a Hausdorff LCVS and $\Lambda$ is countable [i.e., the
  family of seminorms on $X$ is countable], then $X$ is metrizable
with metric
\begin{equation}
\rho(x,y)=\sum^{\infty}_{n=1}2^{-n}\left(\frac{p_{n}(x-y)}{1+p_{n}(x-y)}\right).
\end{equation}
\end{fact}

\begin{definition}
Let $X$ be a locally convex [topological] vector space.
We say $X$ is a \define{Fr\'{e}chet Space} if $X$ is Hausdorff and
$\Lambda$ is countable.
\end{definition}

\begin{example}
The space of continuous real-valued functions $C(\RR^{d})$ or
$C(\Omega)$ when $\Omega$ is an open subset of $\RR^{d}$. For each
compact set $K$ [using the Euclidean topology] in $\Omega$, we define
a seminorm $p_{K} := \|-\|_{C(K)}$ using the supremum norm on $C(K)$
[since $C(K)$ is bounded]. Then $C(\Omega)$ is a Fr\'{e}chet Space
where $\Omega=\bigcup_{n\in\NN}K_{n}$ and $(K_{n})_{n\in\NN}$ allows
us to approximate any $K$ suitably well. Then convergence
$f_{n}\xrightarrow{C(\Omega)}f$ if for every $K$ we have
$f_{n}|_{K}\to F|_{K}$ --- this is called \define{Uniform Convergence
  on Compact Sets}.

We can generalize this to any locally compact Hausdorff space $X$,
which Folland does~\cite[pp.167--168]{folland1999real}.
\end{example}

\begin{example}
The space $C^{\infty}(\RR)$ is a Fr\'{e}chet Space with norm
$\{p_{n,K}\}$ indexed by $n\in\NN$ and $K\subset\RR$ compact, where
$p_{n,K}(f):=\|f^{(n)}\|_{C(K)}$ [writing $f^{(n)}$ for the
  $n^{\text{th}}$ derivative of $f$].
This again is a countable family when we use monotone sequence of $K$
which exaust open sets.
\end{example}

\begin{xca}
The space $C^{\infty}(\RR)$ has the Heine--Borel property (compact
subsets are closed and bounded---``boundedness'' here is with respect
to every seminorm). However, $C(\RR)$ \emph{does not} have the
Heine--Breol property.
\end{xca}

\begin{example}
The space $\ell^{1/2}$ consisting of sequences $(x_{n})$ such that
$\sum_{n}|x_{n}|^{1/2}<\infty$ is finite. Then the function
\begin{equation}
\rho(x,y)=\sum_{n}|x_{n}-y_{n}|^{1/2}
\end{equation}
is a metric on $\ell^{1/2}$.

Note: this is \emph{not} an $\ell^{p}$-space (since the norm would be
$(\sum_{n}|x_{n}|^{1/2})^{2}$ instead). However, this is a topological
vector space with metric $\rho$ inducing its topology. But it is
\emph{not} a locally compact space. (There are no nontrivial convex
open sets.)
\end{example}

\subsection{Weak and Weak-* Topologies}

\begin{node}
Let $X$ be a normed vector space. Let $X'$ be its dual space [so $X'$
  consists of continuous linear functionals].
Then $X'$ is a Banach space. Then by Hanh--Banach, we have 
\begin{equation}
\|x\|:=\sup_{\substack{f\in X'\\ \|f\|\leq1}}|f(x)|,
\end{equation}
taken over the unit ball in $X'$. Compare this to the operator norm
\begin{equation}
\|f\|_{\text{op}}=\sup_{\|x\|\leq1}|f(x)|.
\end{equation}
\end{node}

\begin{definition}
Let $X$ be a normed vector space, let $X'$ be its dual space.
We define the \define{Canonical Embedding} $\pi\colon X\to X''$
by sending $x\mapsto(f\mapsto f(x))$, or $x\mapsto \pi_{x}$ where
$\pi_{x}(f)=f(x)$ for any $f\in X'$.
\end{definition}

\begin{corollary}[of Hanh--Banach]
We have $\|\pi(x)\|=\|x\|$. Therefore $\pi$ is an isometry.
\end{corollary}

\begin{definition}
We say $X$ is \define{Reflexive} if $\pi$ is surjective.
\end{definition}

\begin{example}
Riesz representation theorem implies Hilbert spaces are reflexive.
\end{example}

\begin{node}
Reflexive Banach spaces are ``kind of like'' Hilbert spaces.
\end{node}

\begin{example}
For $1<p<\infty$, $\ell^{p}$-spaces are reflexive.
\end{example}


\begin{example}
We see $(\ell^{2})'\iso\ell^{2}$ canonically isomorphic by Riesz, so
$y\in\ell^{2}\mapsto f_{y}(x)=\sum x_{n}y_{n}$ is standard scalar product.
\end{example}

\begin{example}
We see $(\ell^{1})'\iso\ell^{\infty}$ when we use same formula as
previous example.
\end{example}

\begin{example}
On the other hand, $\ell^{1}\propersubset(\ell^{\infty})'$ but
$\ell^{1}\neq(\ell^{\infty})'$ since $\ell^{1}$ is separable but
$(\ell^{\infty})'$ is not separable.
\end{example}

\begin{example}
The space of continuous functions which converge to zero at infinity
is denoted $C_{0}(-)$. We see $(C_{0}(-))'\iso\ell^{1}$.
Then $C_{0}(-)''\iso\ell^{\infty}$.
\end{example}

\begin{definition}
Let $X$ be a normed vector space.
We define the \define{Weak Topology} on $X$ (denoted (w)) as the
topology generated by $X'$ [i.e., such that $X'$ consists of
  continuous functions].

The weak topology is traditionally denoted $\sigma(X,X')$ where we
understand $\sigma(X,\mathcal{F})$ refers to the weakest topology on $X$
[first slot] such that every function appearing as an element in
$\mathcal{F}\subset\hom(X,Y)$ [the second slot, where $Y$ is some
  topological space] is continuous.

The weak topology is weaker than the normed topology.
\end{definition}

\begin{node}
The space $(X,w)$ with the weak topology is a locally convex space
with seminorms $\{|f|\}_{f\in X'}$.
\end{node}

\begin{node}
In terms of nets, the weak topology can be defined as follows:
$x_{\alpha}\xrightarrow{w}x$ iff $\forall f\in X'\ldotp f(x_{\alpha})\to f(x)$.
\end{node}

\begin{node}
Weak topology is ``just'' the product topology (or the ``pointwise''
topology). 
\end{node}

\begin{remark}
Folland~\cite[pp.168--169]{folland1999real} gives a very expedited
discussion of the weak topology.
\end{remark}